\chapter{Literature Review}
\label{ch:literature}

\section{Introduction}

This chapter reviews the research and systems that Synapse builds upon. It begins
with retrieval-augmented generation, which is the foundation of the grounded
chat, and then looks at the refinements that make retrieval more reliable, such
as self-correction and multi-agent reasoning. It then covers agents that produce
artefacts, knowledge graphs that reveal the structure of a corpus, and the
techniques needed to read documents that are scanned or visually rich. Finally,
it surveys how such systems are evaluated, compares the most relevant systems in
a single table, and identifies the gap that this project fills.

\section{Retrieval-Augmented Generation}

Large language models store a great deal of general knowledge in their
parameters, but that knowledge is fixed at training time and cannot include a
user's private documents. Retrieval-augmented generation, introduced by Lewis et
al.\ \cite{lewis2020rag}, addresses this by attaching a retrieval step to the
model. When a question arrives, the system first finds passages relevant to it
from a trusted collection, and then conditions the model's answer on those
passages. This grounds the answer in real material, allows the source of each
statement to be shown, and lets the knowledge base be updated without retraining
the model.

The quality of the retrieval step is therefore central. Early systems relied on
sparse keyword matching, which is precise for rare terms but misses paraphrases.
Dense passage retrieval \cite{karpukhin2020dpr} instead encodes the question and
the passages as vectors and matches them by meaning, which captures paraphrases
but can drift on exact terms such as names and codes. In practice the two are
complementary, and modern systems combine them. Reciprocal rank fusion
\cite{cormack2009rrf} is a simple and robust way to merge several ranked lists
into one, and is used in Synapse to fuse keyword and vector results. General
purpose text embeddings such as the BGE family \cite{xiao2024cpack} provide the
dense representations, and a cross-encoder reranker re-orders the fused
candidates so that the passages the model finally sees are the most relevant.

\section{Making Retrieval Reliable}

Retrieving relevant passages is necessary but not sufficient. The retrieved
material may still be off-topic, contradictory or simply insufficient, and a
model that answers regardless will produce confident but unsupported text. Two
lines of work address this directly.

Self-RAG \cite{asai2024selfrag} teaches a model to decide when to retrieve, to
criticise its own draft against the retrieved evidence, and to prefer answers
that are well supported. Corrective RAG \cite{yan2024crag} adds a lightweight
evaluator that grades the retrieved set and, when it is weak, triggers a
corrective action such as refining the query or searching more widely before the
answer is composed. Both share a principle that this project adopts: the system
should judge its own evidence and change its behaviour, including declining to
answer, rather than always producing a reply.

A complementary direction breaks the work across several cooperating agents.
MA-RAG \cite{nguyen2025marag} assigns sub-tasks such as planning, retrieval and
synthesis to specialised agents that reason step by step, which helps with
questions that require several pieces of evidence to be combined. Insight-RAG
\cite{pezeshkpour2025insightrag} argues that retrieving on surface similarity
alone can miss the deeper point of a question, and retrieves on the underlying
information need instead. The grounded chat in Synapse follows this overall
shape: it analyses the question, retrieves and grades evidence, and only then
composes a cited answer or abstains.

\section{Agents that Produce Artefacts}

Beyond answering questions, recent agents carry out multi-step tasks. The ReAct
framework \cite{yao2023react} interleaves reasoning and acting, letting a model
think about what to do, take an action, observe the result, and continue. Chain
of thought prompting \cite{wei2022cot} showed that encouraging a model to reason
in explicit steps improves its performance on complex tasks. These ideas underpin
the agent mode in Synapse, which plans its work, decides what data it needs, and
produces an artefact such as a chart or a document.

A practical concern for any agent that reports numbers is faithfulness. If a
model is allowed to write the values inside a chart, it may produce plausible but
incorrect figures. Synapse avoids this by having the agent describe the chart it
wants while the actual data values are bound to it in code from the source file,
so the numbers in an artefact always match the document they came from.

\section{Knowledge Graphs over Documents}

A knowledge graph represents a collection as entities and the relationships
between them, which makes the structure of the collection visible and supports
questions that span several documents. GraphRAG \cite{edge2024graphrag} builds
such a graph from a corpus and uses it to answer broad, summarising questions
that plain passage retrieval handles poorly. CuriousLLM \cite{yang2024curiousllm}
combines graph reasoning with a questioning strategy to improve multi-document
question answering. Synapse includes a knowledge-graph module in the same spirit:
it extracts entities and relationships from a processed library, resolves
duplicates, and presents the result as an interactive graph that keeps a link
from every node and edge back to the passage it came from.

\section{Reading Scanned and Visually Rich Documents}

Real collections are not limited to clean digital text. They include scanned
pages, spreadsheets, slides, figures and tables, and a system that ignores these
will miss much of the content. Document layout analysis locates the regions of a
page, such as paragraphs, tables and figures; DocLayout-YOLO \cite{zhao2024doclayout}
is a recent and efficient model for this task. Optical character recognition then
reads the text inside scanned regions, and modern toolkits such as Surya
\cite{surya2024} handle many languages and layouts. For pages where text alone is
not enough, vision-language models such as Qwen2.5-VL \cite{bai2025qwen25vl} can
read a page image directly and describe its figures and tables. VisDoM
\cite{suri2024visdom} shows the value of treating visually rich elements as
first-class evidence in multi-document question answering. Synapse uses layout
detection, optical character recognition and a vision-language model together so
that scanned and structured documents are processed, not skipped.

\section{Evaluating Document Intelligence}

Measuring these systems is difficult because both retrieval and generation must
be judged, and because a fluent answer is not necessarily a correct one.
Frameworks such as RAGAS \cite{es2024ragas} score properties like faithfulness
and relevance without requiring a hand-written reference for every question.
Because grading free-form answers at scale is costly, a capable model is often
used as an automatic judge \cite{zheng2023judge}, and using more than one judge
helps to expose disagreement. The DOUBLE-BENCH benchmark \cite{shen2025doublebench}
was created specifically to assess document retrieval-augmented generation over
realistic, multimodal documents, with questions that require one or several hops
of reasoning. This project evaluates Synapse on DOUBLE-BENCH together with
purpose-built structured datasets, and reports not only accuracy but also how
often the system is wrongly confident, which is discussed in Chapter~7.

\section{Comparison with Existing Systems}

Table~\ref{tab:comparison} compares Synapse with representative systems across
the capabilities that matter for this project. Most prior systems are strong in
one area, such as corrective retrieval or graph reasoning, but few bring grounded
question answering, artefact generation, knowledge mapping and support for
scanned documents together in a single product that a team can use.

\begin{table}[H]
  \centering
  \caption{Comparison of Synapse with representative related systems. A tick
  means the capability is a designed feature of the system; a dash means it is
  not a focus of that work. Synapse is distinguished by combining grounded,
  abstaining question answering with artefact generation, knowledge mapping and
  support for scanned and structured documents in one platform.}
  \label{tab:comparison}
  \small
  \begin{tabularx}{\textwidth}{Lcccccc}
    \toprule
    \textbf{System} & \textbf{Grounded} & \textbf{Abstains} & \textbf{Multi-} &
    \textbf{Artefact} & \textbf{Knowledge} & \textbf{Scanned/} \\
     & \textbf{QA} & \textbf{when unsure} & \textbf{agent} & \textbf{generation} &
    \textbf{graph} & \textbf{structured} \\
    \midrule
    Classic RAG \cite{lewis2020rag}        & \checkmark & --         & --         & -- & -- & -- \\
    Self-RAG \cite{asai2024selfrag}        & \checkmark & \checkmark & --         & -- & -- & -- \\
    Corrective RAG \cite{yan2024crag}      & \checkmark & \checkmark & --         & -- & -- & -- \\
    MA-RAG \cite{nguyen2025marag}          & \checkmark & \checkmark & \checkmark & -- & -- & -- \\
    GraphRAG \cite{edge2024graphrag}       & \checkmark & --         & --         & -- & \checkmark & -- \\
    VisDoM \cite{suri2024visdom}           & \checkmark & --         & --         & -- & -- & \checkmark \\
    \textbf{Synapse (this work)}           & \checkmark & \checkmark & \checkmark & \checkmark & \checkmark & \checkmark \\
    \bottomrule
  \end{tabularx}
\end{table}

\section{Research Gap}

The review shows steady progress on the individual pieces of document
intelligence, but a gap in how they are brought together. Retrieval has become
more reliable through self-correction and multi-agent reasoning, graphs have made
the structure of a corpus visible, and vision-language models have made scanned
and visually rich pages readable. These advances, however, are usually studied in
isolation and demonstrated on clean benchmarks rather than delivered as a single
product. A user with a real, mixed collection of documents still has no one place
where they can process that collection reliably, ask it questions and receive
honest cited answers, turn it into charts and documents, and explore how its
ideas connect, all while sharing the work with a team.

Synapse addresses this gap. It combines a pipeline that handles scanned and
structured documents, a corrective and multi-agent grounded chat that abstains
when the evidence is insufficient, an agent that produces faithful artefacts, and
an interactive knowledge graph, within one multi-tenant web platform. The
chapters that follow describe its requirements, design, implementation and
evaluation.

\section{Summary}

This chapter reviewed the foundations of the project. It traced retrieval-augmented
generation from its origins to the corrective and multi-agent methods that make
it reliable, surveyed agents that produce artefacts, knowledge graphs over
documents, and the techniques for reading scanned and structured pages, and
looked at how such systems are evaluated. A comparison with representative
systems showed that each addresses part of the problem, and the gap analysis
identified the combined, product-level capability that Synapse provides. The next
chapter sets out the requirements of the system.
